% Purpose: summarize the workflow and the descriptive demonstration in one compact section.
\section{Conclusion}
\label{sec:conclusion}

We have presented the Nisa Radial Control Battery (NRCB), a reproducible
control-validation workflow for radial-structure diagnostics in SPARC-like
rotation-curve data. The workflow is branch-separated: a partial-rank branch
of 134 galaxies and a
spectral-control branch of 76 galaxies are analysed independently and are
not interchangeable. By comparing each observed statistic against multiple
control families, the workflow makes explicit how those families behave. A
label-shuffle comparison and an ordering-preserving exact circular-shift
comparison can lead to different conclusions about the same statistic; the
exact circular shifts act as operational finite-sample controls rather than
physical models. For the amplitude-based spectral-concentration statistic,
equality with amplitude-preserving random-phase controls holds by
construction and is reported as a caveat, and fixed smooth profiles are
deterministic comparators rather than stochastic controls.

We make no physical interpretation, discovery, or model-validation claim;
the demonstration is descriptive and characterizes control behavior on the
two samples. Natural methodological extensions include alternative
detrending, non-cyclic boundary conventions, block or permutation controls,
additional radial templates, and application to broader datasets.
